% --- Archon physics formalization source begin ---
% archon:physics
% archon:covers PhyXMiniProblems/problem_phyx_mini_0254.lean
% archon:source-report reports/phyx_mini/problem_phyx_mini_0254.source.json

\chapter{Physics problem phyx\_mini\_0254}
\label{ch:PhyXMiniProblems_problem_phyx_mini_0254}

\paragraph{Problem source.}
\#\# Physical scenario

Figure a is a partial graph of the position function \$x(t)\$ for a simple harmonic oscillator with an angular frequency of \$1.20\textbackslash{} rad/s\$; Figure b is a partial graph of the corresponding velocity function \$v(t)\$. The vertical axis scales are set by \$x\_s = 5.0\textbackslash{} cm\$ and \$v\_s = 5.0\textbackslash{} cm/s\$.

\#\# Question

What is the phase constant (positive) of the SHM if the position function \$x(t)\$ is in the general form \$x = x\_m \textbackslash{}cos(\textbackslash{}omega t + \textbackslash{}phi)\$?

\#\# Answer choices

- A: A: 0.65 rad

- B: B: 0.86 rad

- C: C: 1.03 rad

- D: D: 1.21 rad

\#\# Figure caption (auxiliary; use the image as primary evidence)

The image consists of two graphs labeled (a) and (b).

1. **Graph (a):**
   - **Axes:**
     - The vertical axis is labeled \textbackslash{}(x\textbackslash{}) (cm).
     - The horizontal axis is labeled \textbackslash{}(t\textbackslash{}).
   - **Plot:**
     - A diagonal line is displayed, starting from a positive \textbackslash{}(x\_s\textbackslash{}) value on the vertical axis and decreasing as \textbackslash{}(t\textbackslash{}) increases, crossing the horizontal axis and extending to a negative \textbackslash{}(-x\_s\textbackslash{}) value.
     - The lines are thick and prominent.
   - **Horizontal Reference Lines:**
     - Horizontal lines are present across the graph, indicating measurement units.
   - **Vertical Reference Lines:**
     - Vertical lines indicate increments of \textbackslash{}(x\_s\textbackslash{}) and \textbackslash{}(-x\_s\textbackslash{}).

2. **Graph (b):**
   - **Axes:**
     - The vertical axis is labeled \textbackslash{}(v\textbackslash{}) (cm/s).
     - The horizontal axis is labeled \textbackslash{}(t\textbackslash{}).
   - **Plot:**
     - A diagonal line is displayed, beginning from a positive \textbackslash{}(v\_s\textbackslash{}) value on the vertical axis and decreasing as \textbackslash{}(t\textbackslash{}) increases, crossing the horizontal axis and extending to a negative \textbackslash{}(-v\_s\textbackslash{}) value.
     - This line is also thick and prominent.
   - **Horizontal Reference Lines:**
     - Parallel horizontal lines are present along the graph’s height.
   - **Vertical Reference Lines:**
     - Vertical lines indicate increments of \textbackslash{}(v\_s\textbackslash{}) and \textbackslash{}(-v\_s\textbackslash{}).

Both graphs illustrate a linear relationship with time (\textbackslash{}(t\textbackslash{})), each capturing a negative slope line, and include labeled increments to suggest specific values or scales.

\#\# Dataset metadata

Resonance and Harmonics; Physical Model Grounding Reasoning, Spatial Relation Reasoning

\paragraph{Recorded answer/context.}
C: C: 1.03 rad

\paragraph{Figure/image path.}
/root/proposal\_for\_physic/science-mango/phyx\_mini\_run/phyx\_data/test\_image/254.png

\paragraph{Formalization target.}
create a compiling Lean file with sorry bodies at `PhyXMiniProblems/problem\_phyx\_mini\_0254.lean`. The Lean declarations must preserve the physical quantities, dimensions or dimensional roles, figure labels, governing-law hypotheses, and final relation expressed by this problem.
Use Mathlib/Physlib names found through LeanExplore where available. If a physics API is missing, introduce faithful local abstractions rather than scalar placeholder aliases.

\begin{theorem}[Physics formalization target]
\label{thm:physics:phyx_mini_0254:target}
\lean{PhyXMiniProblems.ProblemPhyXMini0254.phaseConstantFromPositionAndVelocityGraphs}
\uses{def:physics:phyx-mini-0254:phyxminiproblems-problemphyxmini0254-harmonicphasesetup, def:physics:phyx-mini-0254:phyxminiproblems-problemphyxmini0254-matchesstatedangularfrequency, def:physics:phyx-mini-0254:phyxminiproblems-problemphyxmini0254-matchessuppliedphasegraphs, def:physics:phyx-mini-0254:phyxminiproblems-problemphyxmini0254-hasphysicalharmonicparameters, def:physics:phyx-mini-0254:phyxminiproblems-problemphyxmini0254-satisfiespositivephasecosineshm, def:physics:phyx-mini-0254:phyxminiproblems-problemphyxmini0254-answerchoice, def:physics:phyx-mini-0254:phyxminiproblems-problemphyxmini0254-answerchoiceinradians, def:physics:phyx-mini-0254:phyxminiproblems-problemphyxmini0254-recordedanswerchoice, def:physics:phyx-mini-0254:phyxminiproblems-problemphyxmini0254-matchesanswertonearesthundredth}
The assigned autoformalize agent should translate this physics problem into Lean declarations in the covered file, with theorem and lemma proof bodies written as `by sorry`.
\end{theorem}
\begin{proof}
This is an autoformalization task, not a proof task. Produce statements that can later be proved by the physics prover without weakening the physical model.
\end{proof}

% --- Archon declaration topology begin ---
\section{Declaration topology}
\begin{definition}[Time Quantity]
  \label{def:physics:phyx-mini-0254:phyxminiproblems-problemphyxmini0254-timequantity}
  \lean{PhyXMiniProblems.ProblemPhyXMini0254.TimeQuantity}
  A unit-independent physical time
\end{definition}
\begin{proof}
  This is the declared model definition of Time Quantity.
\end{proof}
\begin{definition}[Length Quantity]
  \label{def:physics:phyx-mini-0254:phyxminiproblems-problemphyxmini0254-lengthquantity}
  \lean{PhyXMiniProblems.ProblemPhyXMini0254.LengthQuantity}
  A unit-independent signed position or displacement
\end{definition}
\begin{proof}
  This is the declared model definition of Length Quantity.
\end{proof}
\begin{definition}[Speed Quantity]
  \label{def:physics:phyx-mini-0254:phyxminiproblems-problemphyxmini0254-speedquantity}
  \lean{PhyXMiniProblems.ProblemPhyXMini0254.SpeedQuantity}
  A unit-independent signed velocity component
\end{definition}
\begin{proof}
  This is the declared model definition of Speed Quantity.
\end{proof}
\begin{definition}[Angular Frequency Quantity]
  \label{def:physics:phyx-mini-0254:phyxminiproblems-problemphyxmini0254-angularfrequencyquantity}
  \lean{PhyXMiniProblems.ProblemPhyXMini0254.AngularFrequencyQuantity}
  Angular frequency, with inverse-time dimension; radians are dimensionless
\end{definition}
\begin{proof}
  This is the declared model definition of Angular Frequency Quantity.
\end{proof}
\begin{definition}[time In Seconds]
  \label{def:physics:phyx-mini-0254:phyxminiproblems-problemphyxmini0254-timeinseconds}
  \lean{PhyXMiniProblems.ProblemPhyXMini0254.timeInSeconds}
  \uses{def:physics:phyx-mini-0254:phyxminiproblems-problemphyxmini0254-timequantity}
  Read a physical time in seconds
\end{definition}
\begin{proof}
  This is the declared model definition of time In Seconds.
\end{proof}
\begin{definition}[length In Centimeters]
  \label{def:physics:phyx-mini-0254:phyxminiproblems-problemphyxmini0254-lengthincentimeters}
  \lean{PhyXMiniProblems.ProblemPhyXMini0254.lengthInCentimeters}
  \uses{def:physics:phyx-mini-0254:phyxminiproblems-problemphyxmini0254-lengthquantity}
  Read a signed physical length in centimeters
\end{definition}
\begin{proof}
  This is the declared model definition of length In Centimeters.
\end{proof}
\begin{definition}[speed In Centimeters Per Second]
  \label{def:physics:phyx-mini-0254:phyxminiproblems-problemphyxmini0254-speedincentimeterspersecond}
  \lean{PhyXMiniProblems.ProblemPhyXMini0254.speedInCentimetersPerSecond}
  \uses{def:physics:phyx-mini-0254:phyxminiproblems-problemphyxmini0254-speedquantity}
  Read a signed physical speed in centimeters per second
\end{definition}
\begin{proof}
  This is the declared model definition of speed In Centimeters Per Second.
\end{proof}
\begin{definition}[angular Frequency In Radians Per Second]
  \label{def:physics:phyx-mini-0254:phyxminiproblems-problemphyxmini0254-angularfrequencyinradianspersecond}
  \lean{PhyXMiniProblems.ProblemPhyXMini0254.angularFrequencyInRadiansPerSecond}
  \uses{def:physics:phyx-mini-0254:phyxminiproblems-problemphyxmini0254-angularfrequencyquantity}
  Read an angular frequency in radians per second
\end{definition}
\begin{proof}
  This is the declared model definition of angular Frequency In Radians Per Second.
\end{proof}
\begin{definition}[Figure Panel]
  \label{def:physics:phyx-mini-0254:phyxminiproblems-problemphyxmini0254-figurepanel}
  \lean{PhyXMiniProblems.ProblemPhyXMini0254.FigurePanel}
  The two panels explicitly labeled in the supplied image
\end{definition}
\begin{proof}
  This is the declared model definition of Figure Panel.
\end{proof}
\begin{definition}[Axis Quantity]
  \label{def:physics:phyx-mini-0254:phyxminiproblems-problemphyxmini0254-axisquantity}
  \lean{PhyXMiniProblems.ProblemPhyXMini0254.AxisQuantity}
  Quantities printed on the horizontal and vertical graph axes
\end{definition}
\begin{proof}
  This is the declared model definition of Axis Quantity.
\end{proof}
\begin{definition}[Vertical Scale Label]
  \label{def:physics:phyx-mini-0254:phyxminiproblems-problemphyxmini0254-verticalscalelabel}
  \lean{PhyXMiniProblems.ProblemPhyXMini0254.VerticalScaleLabel}
  The scale symbols printed next to the vertical axes
\end{definition}
\begin{proof}
  This is the declared model definition of Vertical Scale Label.
\end{proof}
\begin{definition}[Phase Graph Figure]
  \label{def:physics:phyx-mini-0254:phyxminiproblems-problemphyxmini0254-phasegraphfigure}
  \lean{PhyXMiniProblems.ProblemPhyXMini0254.PhaseGraphFigure}
  \uses{def:physics:phyx-mini-0254:phyxminiproblems-problemphyxmini0254-timequantity, def:physics:phyx-mini-0254:phyxminiproblems-problemphyxmini0254-lengthquantity, def:physics:phyx-mini-0254:phyxminiproblems-problemphyxmini0254-speedquantity, def:physics:phyx-mini-0254:phyxminiproblems-problemphyxmini0254-figurepanel, def:physics:phyx-mini-0254:phyxminiproblems-problemphyxmini0254-axisquantity, def:physics:phyx-mini-0254:phyxminiproblems-problemphyxmini0254-verticalscalelabel}
  The labels, scales, common origin, and qualitative slope information read from the two graph panels. The scale values and curve values are independent physical quantities; their numerical relations are recorded separately in MatchesSuppliedPhaseGraphs
\end{definition}
\begin{proof}
  This is the declared model definition of Phase Graph Figure.
\end{proof}
\begin{definition}[Harmonic Phase Setup]
  \label{def:physics:phyx-mini-0254:phyxminiproblems-problemphyxmini0254-harmonicphasesetup}
  \lean{PhyXMiniProblems.ProblemPhyXMini0254.HarmonicPhaseSetup}
  \uses{def:physics:phyx-mini-0254:phyxminiproblems-problemphyxmini0254-timequantity, def:physics:phyx-mini-0254:phyxminiproblems-problemphyxmini0254-lengthquantity, def:physics:phyx-mini-0254:phyxminiproblems-problemphyxmini0254-speedquantity, def:physics:phyx-mini-0254:phyxminiproblems-problemphyxmini0254-angularfrequencyquantity, def:physics:phyx-mini-0254:phyxminiproblems-problemphyxmini0254-phasegraphfigure}
  The oscillator quantities named in the problem. In particular, the phase constant is an independent dimensionless parameter, not a definition of the requested numerical answer
\end{definition}
\begin{proof}
  This is the declared model definition of Harmonic Phase Setup.
\end{proof}
\begin{definition}[Matches Stated Angular Frequency]
  \label{def:physics:phyx-mini-0254:phyxminiproblems-problemphyxmini0254-matchesstatedangularfrequency}
  \lean{PhyXMiniProblems.ProblemPhyXMini0254.MatchesStatedAngularFrequency}
  \uses{def:physics:phyx-mini-0254:phyxminiproblems-problemphyxmini0254-angularfrequencyinradianspersecond, def:physics:phyx-mini-0254:phyxminiproblems-problemphyxmini0254-harmonicphasesetup}
  The stated angular frequency 1.20 rad/s
\end{definition}
\begin{proof}
  This is the declared model definition of Matches Stated Angular Frequency.
\end{proof}
\begin{definition}[Matches Supplied Phase Graphs]
  \label{def:physics:phyx-mini-0254:phyxminiproblems-problemphyxmini0254-matchessuppliedphasegraphs}
  \lean{PhyXMiniProblems.ProblemPhyXMini0254.MatchesSuppliedPhaseGraphs}
  \uses{def:physics:phyx-mini-0254:phyxminiproblems-problemphyxmini0254-timeinseconds, def:physics:phyx-mini-0254:phyxminiproblems-problemphyxmini0254-lengthincentimeters, def:physics:phyx-mini-0254:phyxminiproblems-problemphyxmini0254-speedincentimeterspersecond, def:physics:phyx-mini-0254:phyxminiproblems-problemphyxmini0254-harmonicphasesetup}
  Labels and quantitative readouts from the primary image. The top and bottom scale labels are separated from zero by five equal graph divisions. At the common vertical axis the position curve is two divisions above zero, whereas the velocity curve is four divisions below zero. Thus the origin values are respectively 2/5 of x\_s and -4/5 of v\_s. These are figure observations, not assumptions about the requested phase
\end{definition}
\begin{proof}
  This is the declared model definition of Matches Supplied Phase Graphs.
\end{proof}
\begin{definition}[Has Physical Harmonic Parameters]
  \label{def:physics:phyx-mini-0254:phyxminiproblems-problemphyxmini0254-hasphysicalharmonicparameters}
  \lean{PhyXMiniProblems.ProblemPhyXMini0254.HasPhysicalHarmonicParameters}
  \uses{def:physics:phyx-mini-0254:phyxminiproblems-problemphyxmini0254-lengthincentimeters, def:physics:phyx-mini-0254:phyxminiproblems-problemphyxmini0254-angularfrequencyinradianspersecond, def:physics:phyx-mini-0254:phyxminiproblems-problemphyxmini0254-harmonicphasesetup}
  Positivity and the conventional positive representative for the phase. Restricting φ to one positive turn selects a representative of the periodic cosine phase but does not prescribe its value
\end{definition}
\begin{proof}
  This is the declared model definition of Has Physical Harmonic Parameters.
\end{proof}
\begin{definition}[Satisfies Positive Phase Cosine SHM]
  \label{def:physics:phyx-mini-0254:phyxminiproblems-problemphyxmini0254-satisfiespositivephasecosineshm}
  \lean{PhyXMiniProblems.ProblemPhyXMini0254.SatisfiesPositivePhaseCosineSHM}
  \uses{def:physics:phyx-mini-0254:phyxminiproblems-problemphyxmini0254-timequantity, def:physics:phyx-mini-0254:phyxminiproblems-problemphyxmini0254-timeinseconds, def:physics:phyx-mini-0254:phyxminiproblems-problemphyxmini0254-lengthincentimeters, def:physics:phyx-mini-0254:phyxminiproblems-problemphyxmini0254-speedincentimeterspersecond, def:physics:phyx-mini-0254:phyxminiproblems-problemphyxmini0254-angularfrequencyinradianspersecond, def:physics:phyx-mini-0254:phyxminiproblems-problemphyxmini0254-harmonicphasesetup}
  The governing simple-harmonic-motion laws in the sign convention stated in the question: x(t) = x\_m cos (ωt + φ) and v(t) = -ω x\_m sin (ωt + φ). The second field states that the plotted velocity is the velocity corresponding to the plotted position. Neither field mentions an answer choice or a numerical phase
\end{definition}
\begin{proof}
  This is the declared model definition of Satisfies Positive Phase Cosine SHM.
\end{proof}
\begin{definition}[Answer Choice]
  \label{def:physics:phyx-mini-0254:phyxminiproblems-problemphyxmini0254-answerchoice}
  \lean{PhyXMiniProblems.ProblemPhyXMini0254.AnswerChoice}
  Labels of the four displayed phase-constant choices
\end{definition}
\begin{proof}
  This is the declared model definition of Answer Choice.
\end{proof}
\begin{definition}[answer Choice In Radians]
  \label{def:physics:phyx-mini-0254:phyxminiproblems-problemphyxmini0254-answerchoiceinradians}
  \lean{PhyXMiniProblems.ProblemPhyXMini0254.answerChoiceInRadians}
  \uses{def:physics:phyx-mini-0254:phyxminiproblems-problemphyxmini0254-answerchoice}
  The phase in radians printed beside each answer label
\end{definition}
\begin{proof}
  This is the declared model definition of answer Choice In Radians.
\end{proof}
\begin{definition}[recorded Answer Choice]
  \label{def:physics:phyx-mini-0254:phyxminiproblems-problemphyxmini0254-recordedanswerchoice}
  \lean{PhyXMiniProblems.ProblemPhyXMini0254.recordedAnswerChoice}
  \uses{def:physics:phyx-mini-0254:phyxminiproblems-problemphyxmini0254-answerchoice}
  The answer label recorded in the supplied dataset metadata
\end{definition}
\begin{proof}
  This is the declared model definition of recorded Answer Choice.
\end{proof}
\begin{definition}[Matches Answer To Nearest Hundredth]
  \label{def:physics:phyx-mini-0254:phyxminiproblems-problemphyxmini0254-matchesanswertonearesthundredth}
  \lean{PhyXMiniProblems.ProblemPhyXMini0254.MatchesAnswerToNearestHundredth}
  \uses{def:physics:phyx-mini-0254:phyxminiproblems-problemphyxmini0254-answerchoice, def:physics:phyx-mini-0254:phyxminiproblems-problemphyxmini0254-answerchoiceinradians}
  Agreement with a displayed phase to the nearest 0.01 rad
\end{definition}
\begin{proof}
  This is the declared model definition of Matches Answer To Nearest Hundredth.
\end{proof}
% --- Archon declaration topology end ---
% --- Archon physics formalization source end ---
